% =====================================================================
%  preamble.tex -- packages, theorem environments, epistemic status tags
% =====================================================================

\usepackage[margin=1.1in]{geometry}
\usepackage{amsmath,amssymb,amsthm,mathtools}
\usepackage{enumitem}
\usepackage{xcolor}
\usepackage{mdframed}
\usepackage{tikz}
\usetikzlibrary{arrows.meta,positioning,calc,decorations.pathmorphing,fit,backgrounds}
\usepackage{booktabs}
\usepackage{array}      % >{...} column specifiers, used for the tag tables
\usepackage{longtable}
\usepackage{caption}
\usepackage{fancyhdr}
\usepackage{titlesec}
\usepackage[colorlinks=true,linkcolor=linkblue,citecolor=linkblue,
            urlcolor=linkblue]{hyperref}

% ---------------------------------------------------------------------
%  Colours
% ---------------------------------------------------------------------
\definecolor{linkblue}{RGB}{28,78,140}
\definecolor{intuitioncol}{RGB}{38,92,158}
\definecolor{honestycol}{RGB}{168,52,40}
\definecolor{examplecol}{RGB}{34,110,72}
\definecolor{warncol}{RGB}{176,96,16}
\definecolor{greybg}{RGB}{246,246,244}

% ---------------------------------------------------------------------
%  Theorem environments -- numbered within chapters
% ---------------------------------------------------------------------
\theoremstyle{definition}
\newtheorem{definition}{Definition}[chapter]
\newtheorem{construction}[definition]{Construction}
\newtheorem{example}[definition]{Example}
\newtheorem{exercise}[definition]{Exercise}
\newtheorem{notation}[definition]{Notation}

\theoremstyle{plain}
\newtheorem{theorem}[definition]{Theorem}
\newtheorem{proposition}[definition]{Proposition}
\newtheorem{lemma}[definition]{Lemma}
\newtheorem{corollary}[definition]{Corollary}

\theoremstyle{remark}
\newtheorem{remark}[definition]{Remark}
\newtheorem{openproblem}[definition]{Open Problem}

% ---------------------------------------------------------------------
%  Epistemic status tags
%
%  Every substantive claim in this book carries one of these.  The point
%  is that the source documents disagree with each other, and a smooth
%  narrative would hide that.  Read the tag before you trust the claim.
% ---------------------------------------------------------------------
\newcommand{\statusbadge}[2]{%
  \textcolor{#1}{\small\textsf{\textbf{[#2]}}}}

\newcommand{\Proved}{\statusbadge{examplecol}{PROVED}}
\newcommand{\ProvedCond}{\statusbadge{examplecol}{PROVED UNDER STATED CONDITIONS}}
\newcommand{\Implemented}{\statusbadge{warncol}{IMPLEMENTED, UNPROVED}}
\newcommand{\Proposed}{\statusbadge{linkblue}{PROPOSED}}
\newcommand{\Refuted}{\statusbadge{honestycol}{REFUTED}}
\newcommand{\Contested}{\statusbadge{honestycol}{CONTESTED BY AUDIT}}

% Measured: a number was printed by a program that can be re-run.  Distinct
% from PROVED (an argument) and from IMPLEMENTED (code exists, no result
% quoted).
\newcommand{\Measured}{\statusbadge{examplecol}{MEASURED}}

% ---------------------------------------------------------------------
%  Disposition tags
%
%  Used by chapters that are still being brought up to the standard of
%  Parts I and II.  A chapter carrying one is a chapter whose framing is
%  older than the foundations it now rests on.
% ---------------------------------------------------------------------
\newcommand{\dispbadge}[2]{%
  \textcolor{#1}{\small\textsf{\textbf{<#2>}}}}

\newcommand{\Keep}{\dispbadge{examplecol}{KEEP}}
\newcommand{\Recast}{\dispbadge{warncol}{RECAST}}
\newcommand{\Retire}{\dispbadge{honestycol}{RETIRE}}
\newcommand{\NewV}{\dispbadge{linkblue}{NEW}}

% ---------------------------------------------------------------------
%  Layered-exposition boxes
% ---------------------------------------------------------------------
\newmdenv[topline=false,bottomline=false,rightline=false,linewidth=2pt,
  linecolor=intuitioncol,innerleftmargin=10pt,innertopmargin=3pt,
  innerbottommargin=3pt,skipabove=8pt,skipbelow=8pt]{intu}
\newmdenv[topline=false,bottomline=false,rightline=false,linewidth=2pt,
  linecolor=honestycol,innerleftmargin=10pt,innertopmargin=3pt,
  innerbottommargin=3pt,skipabove=8pt,skipbelow=8pt]{hon}
\newmdenv[topline=false,bottomline=false,rightline=false,linewidth=2pt,
  linecolor=examplecol,innerleftmargin=10pt,innertopmargin=3pt,
  innerbottommargin=3pt,skipabove=8pt,skipbelow=8pt]{workedbox}
\newmdenv[linewidth=0.6pt,linecolor=black!35,backgroundcolor=greybg,
  innerleftmargin=9pt,innerrightmargin=9pt,innertopmargin=7pt,
  innerbottommargin=7pt,skipabove=9pt,skipbelow=9pt,roundcorner=2pt]{sourcebox}

\newcommand{\Intu}{\noindent\textbf{\color{intuitioncol}Intuition.}\ }
\newcommand{\Hon}{\noindent\textbf{\color{honestycol}Claims and non-claims.}\ }
\newcommand{\Worked}{\noindent\textbf{\color{examplecol}Worked example.}\ }

% Provenance note: which file in DOCS/ a passage descends from.
\newcommand{\source}[1]{%
  \begin{sourcebox}\small\textsf{\textbf{Source.}} #1\end{sourcebox}}

% ---------------------------------------------------------------------
%  File and path names
%
%  \texttt{DOCS/SOME_LONG_NAME.md} does not hyphenate and contains no
%  legal breakpoint, so long paths ran off the right margin -- they were
%  the largest single source of overfull lines in the book.  \file{}
%  breaks at / and _ instead.  Underscores inside \file{} are LITERAL:
%  write \file{DOCS/MATH_TOOLBOX.md}, not \file{DOCS/MATH\_TOOLBOX.md}.
%  (\DeclareUrlCommand comes from url, already loaded by hyperref.)
% ---------------------------------------------------------------------
\DeclareUrlCommand\file{\urlstyle{tt}}

% ---------------------------------------------------------------------
%  Work plans
%
%  The rewrite is done section by section, and each section states its
%  own plan before it is written.  A work plan is not commentary on the
%  finished text -- it is the instruction for producing it, and it stays
%  in the document until the section it governs is complete, at which
%  point it is replaced by the prose it asked for.
%
%  Usage:
%    \begin{workplan}{<disposition tag>}
%      \wpgoal{one sentence: what this section must establish}
%      \wpsource{where the raw material is}
%      \wpsteps{\item ... \item ...}
%      \wpdone{the condition under which this section is finished}
%    \end{workplan}
% ---------------------------------------------------------------------
\newmdenv[linewidth=1pt,linecolor=linkblue!60,backgroundcolor=linkblue!4,
  innerleftmargin=10pt,innerrightmargin=10pt,innertopmargin=8pt,
  innerbottommargin=8pt,skipabove=10pt,skipbelow=10pt,roundcorner=3pt]{wpbox}

\newenvironment{workplan}[1]
  {\begin{wpbox}\small\noindent\textsf{\textbf{\color{linkblue}Work plan.}}\ #1\par\vspace{3pt}}
  {\end{wpbox}}

\newcommand{\wpgoal}[1]{\noindent\textit{Goal.}\ #1\par\vspace{2pt}}
\newcommand{\wpsource}[1]{\noindent\textit{Raw material.}\ #1\par\vspace{2pt}}
\newcommand{\wpblocked}[1]{\noindent\textit{Blocked by.}\ #1\par\vspace{2pt}}
\newcommand{\wpsteps}[1]{\noindent\textit{Steps.}\begin{enumerate}[leftmargin=1.4em,itemsep=1pt]#1\end{enumerate}}
\newcommand{\wpdone}[1]{\vspace{2pt}\noindent\textit{Done when.}\ #1\par}

% ---------------------------------------------------------------------
%  Section agendas
%
%  A work plan is an instruction for writing a section that does not yet
%  exist.  An agenda is different: it is a standing statement of where a
%  *written* section actually stands -- what it has settled, what it has
%  not, and the ordered steps that would close the gap.  Every section of
%  Parts I and II carries one, and it stays there until its open list is
%  empty.  The two must not be confused: a work plan being present means
%  no prose exists; an agenda being present means prose exists and is
%  incomplete in stated ways.
%
%  Usage:
%    \begin{agenda}
%      \agset{\item ...}      % established, here or by citation
%      \agopen{\item ...}     % not established; the honest gaps
%      \agsteps{\item ...}    % ordered work that would close them
%    \end{agenda}
% ---------------------------------------------------------------------
\newmdenv[linewidth=1pt,linecolor=examplecol!55,backgroundcolor=examplecol!5,
  innerleftmargin=10pt,innerrightmargin=10pt,innertopmargin=8pt,
  innerbottommargin=8pt,skipabove=10pt,skipbelow=10pt,roundcorner=3pt]{agbox}

\newenvironment{agenda}
  {\begin{agbox}\small\noindent\textsf{\textbf{\color{examplecol}Agenda.}}\par\vspace{2pt}}
  {\end{agbox}}

\newcommand{\agset}[1]{\noindent\textit{Settled.}%
  \begin{enumerate}[leftmargin=1.6em,itemsep=1pt,topsep=2pt,label=\textup{S\arabic*.}]#1\end{enumerate}\vspace{3pt}}
\newcommand{\agopen}[1]{\noindent\textit{Open.}%
  \begin{enumerate}[leftmargin=1.6em,itemsep=1pt,topsep=2pt,label=\textup{O\arabic*.}]#1\end{enumerate}\vspace{3pt}}
\newcommand{\agsteps}[1]{\noindent\textit{Next steps, in order.}%
  \begin{enumerate}[leftmargin=1.6em,itemsep=1pt,topsep=2pt,label=\textup{\arabic*.}]#1\end{enumerate}}

% ---------------------------------------------------------------------
%  Headers
% ---------------------------------------------------------------------
\pagestyle{fancy}
\fancyhf{}
\fancyhead[LE,RO]{\small\thepage}
\fancyhead[RE]{\small\itshape\nouppercase{\leftmark}}
\fancyhead[LO]{\small\itshape\nouppercase{\rightmark}}
\renewcommand{\headrulewidth}{0.4pt}

\setlist{itemsep=2pt,parsep=0pt,topsep=4pt}
\allowdisplaybreaks
